\justifying
\NSFCBodyText

\subsubsection{与本项目相关的研究工作积累}

申请人为新疆大学副教授、硕士研究生导师，研究积累涉及多尺度时序建模、多源信息融合与视觉问答。与本项目最直接的联系在于非平稳数据表示、不同来源证据的组织和可复现实验设计。上述积累可以支持动态任务构造、轨迹证据表示与方法比较；本项目的智能体持续后训练、适配器记忆及行动校准仍需专门研究和验证。

在非平稳时序建模方面，申请人参与完成 DWT-Former 与 WaveKAN 研究。DWT-Former 融合小波多尺度特征与 Transformer 时间表示，开展光伏功率预测\cite{dwtformer}。

WaveKAN 研究自适应小波分解与 Kolmogorov--Arnold 网络的长时序建模\cite{wavekan}。这两项工作支撑对尺度变化、时间划分和预测误差的理解，可用于构造具有明确时间边界的动态补充任务。项目将把任务层的工具调用与底层预测器效果分开评价，避免把预测模型进步直接算为 Agent 方法收益。

在证据表示方面，申请人开展了知识感知视觉问答研究，通过多层视觉表示增强图像理解\cite{knowledgevqa}。该工作支撑对异质证据、问题条件和模型判断之间关系的分析，可迁移为本项目的轨迹证据组织思路。此次研究以文本工具任务为主体，已有多模态成果只用于说明相关方法能力，不扩展出缺少验证资源的新研究路线。

\subsubsection{已取得的研究工作成绩}

申请人为 DWT-Former 的通讯作者，该成果发表于 \textit{Energy}，DOI 为 10.1016/j.energy.2025.139283；为 WaveKAN 的共同第一作者，该成果发表于 \textit{Neurocomputing}，DOI 为 10.1016/j.neucom.2026.133361；为 Knowledge-aware Image Understanding with Multi-level Visual Representation Enhancement for Visual Question Answering 的第一作者兼通讯作者，该成果发表于 \textit{Machine Learning}，DOI 为 10.1007/s10994-023-06426-6。作者身份依据申请人已确认资料，题录依据出版登记元数据核对。

申请人主持多模态技术、光伏功率预测及知识库问答相关项目，并参与高效注意力机制研究，具备组织模型研究和软件实现的经验。有关项目的正式名称、经费、执行期及本项目边界在承担项目部分列示。已有成果能够说明研究积累，但不作为 PACT-Agent 已实现性能提升的证据；本项目将新建专用工具沙箱、轨迹库和校准实验，所有先导结果均以实际运行记录为准。

\subsubsection{可行性分析}

\textbf{理论与方法方面，}参数高效训练、经验回放和概率校准均具有明确基础方法，能够先形成独立基线再开展机制研究。任务证据是否有效、轨迹选择是否放大验证偏差以及校准是否随策略更新失效，都可以通过受控干预与组合对照检验，不依赖无法观察的业务结果。

\textbf{实验与实现方面，}已有多尺度时序和问答研究为任务构造、误差分析及模型实现提供经验；Linux、Python、深度学习框架、Git 与 LaTeX 工具链已具备。研究从可重置的文本工具环境开始，把外部接口变化、随机用户模拟和模型解码分别控制，先回答可识别的问题，再扩大任务范围。

\textbf{组织方面，}由申请人负责总体研究、关键方法设计与质量把关。模型训练、轨迹审计、持续评估和工程整合的具体承担人员为\textbf{【待补 F-PACT-02：本项目成员及确认分工】}。现有其他项目的学生名单不自动作为本项目团队；相关人力投入需在首阶段落实。任务之间采用明确接口，人员未完全到位时优先完成单模型、单沙箱的核心实验。

\subsubsection{研究风险与应对措施}

\textbf{证据偏差与误差放大风险。}生成器和验证器可能对同类错误同时自信，使错误轨迹在筛选与重放中反复出现。早期信号是执行判定与验证器高分持续不一致，或加入高分轨迹后独立开发集反而退化。预案是增加确定性约束检查与分层人工审计，将无法判定轨迹保留为未知；暂停其进入偏好对和历史记忆，用仅可靠结果过滤的基线维持研究。相关错误干预本身保留为主要失败边界实验。

\textbf{方法增益不足与任务干扰风险。}动态分配和双重记忆可能增加开销，却未带来等预算收益。早期信号是总GPU时匹配后优势消失，或适配器合并引起历史开发集持续下降。预案是保留静态配置、固定容量回放和独立校准的最小系统，通过模块与交互消融确定无效环节；不以增加模型规模弥补缺少机制证据。最终可交付资源匹配下的边界结论。

\textbf{校准失效与标签不足风险。}工具突变或反馈延迟可能使当前校准器失效。若连续监测窗口中错误执行增多或有效标签不足，则冻结新的高风险状态修改、限制重规划次数并记录覆盖率损失；在沙箱补齐独立校准数据后重新评估。仅报告具备样本支持的风险区间，不能由长期平均覆盖推断当前动作已安全。

\textbf{资源与进度风险。}多种顺序、规模和控制器的组合可能超出首年算力与人力。以首阶段实测成本作为预警依据，若关键复现未按阶段完成，优先保留3B级模型、核心工具任务和三类单独变化，减少外部环境与大模型扩展。飞桨兼容工作先验证核心低秩训练和日志接口，不承诺未经测试的全栈兼容；相关技术限制连同可复现基线一并交付。
